% !TEX TS-program = pdflatex
% !TEX encoding = UTF-8 Unicode
% Couette flow: a page of the wiki "Enhanced dissipation for passive scalars", generated from the wiki on 11 November 2026.
% It compiles as it is with pdflatex (Overleaf, TeXShop, TeXstudio, ...); typeset twice for the citations.
% To propose a change: keep the preamble and the line \title{...}; keep only the blocks you change
% (a \section{...} with its content, or a statement or proof with its \label{...}), change them, and post
% the file in a new topic of the Zulip channel of the part. The rules: page "How to contribute" of the wiki.
\documentclass[11pt]{article}
\usepackage[utf8]{inputenc}
\usepackage[T1]{fontenc}
\usepackage{lmodern}
\usepackage[a4paper,margin=2.5cm]{geometry}
\usepackage{amsmath,amssymb,amsthm}
\usepackage{graphicx,tabularx,listings}
\usepackage[skins,breakable]{tcolorbox}
\usepackage[round]{natbib}
\usepackage[hidelinks]{hyperref}
\setlength{\emergencystretch}{2em}

\newtheorem{theorem}{Theorem}
\newtheorem{proposition}[theorem]{Proposition}
\newtheorem{lemma}[theorem]{Lemma}
\newtheorem{corollary}[theorem]{Corollary}
\newtheorem{conjecture}[theorem]{Conjecture}
\newtheorem*{claim}{Claim}
\theoremstyle{definition}
\newtheorem{definition}[theorem]{Definition}
\newtheorem{example}[theorem]{Example}
\theoremstyle{remark}
\newtheorem{remark}[theorem]{Remark}

%% The commands of the wiki. This preamble is the same for every page: do not change it, and do not
%% define commands of your own (the wiki has none).
\makeatletter
% after a statement: \begin{idea}, \begin{architecture}, \begin{note}[Title] (and \begin{proof}[Title])
\newcommand{\ed@runin}[1]{\par\topsep6\p@\@plus6\p@\relax\trivlist\item[\hskip\labelsep\itshape#1\@addpunct{.}]\ignorespaces}
\newenvironment{idea}[1][Idea]{\ed@runin{#1}}{\endtrivlist\@endpefalse}
\newenvironment{architecture}[1][Proof architecture]{\ed@runin{#1}}{\endtrivlist\@endpefalse}
% \begin{note}[Title]: an aside, in a dashed box; its title in italics at the start of the text
\tcbset{ed@note/.style={enhanced jigsaw,breakable,frame hidden,interior hidden,arc=2pt,borderline={0.5pt}{0pt}{black!50,dash pattern=on 3pt off 2pt},
  left=3mm,right=3mm,top=2mm,bottom=2mm,before skip=3mm plus 1mm,after skip=3mm plus 1mm}}
\newenvironment{note}[1][Note]{\begin{tcolorbox}[ed@note]\noindent{\itshape#1\@addpunct{.}}\enspace\ignorespaces}{\end{tcolorbox}}
% text still to be written; labels of the text; the Lean banner of the statement above
\newcommand{\placeholder}[1]{\if@inlabel\else\par\noindent\fi\textit{[#1]}\@ifnextchar\end{}{\par}}
\newcommand{\ed@label}[1]{\mbox{\footnotesize[#1]}}
\newcommand{\demo}[1]{\ed@label{demo text}}
\newcommand{\heuristic}[1]{\ed@label{heuristic}}
\newcommand{\numerics}[1]{\ed@label{numerics}}
\newcommand{\published}[1]{\ed@label{published}}
\newcommand{\preprint}[1]{\ed@label{preprint, to be verified}}
\newcommand{\raw}[1]{\ed@label{raw, unchecked}}
\def\ed@formalized{formalized}\def\ed@image{image}\def\ed@video{video}\def\ed@diagram{diagram}
\newcommand{\lean}[2][]{\par{\noindent\footnotesize\def\ed@a{#2}Lean: \ifx\ed@a\ed@formalized formalized\ifx\relax#1\relax\else\
  (\href{#1}{certificate})\fi\else not yet formalized\fi.\par}}
% the contributors of a section, most recent first (chosen by the editors), in small type under its title
\newcommand{\contributors}[1]{\par{\noindent\footnotesize\itshape Contributors: #1\par}\smallskip}
% the two lines under the title: the page's team and status, its Zulip topic (default: the title)
\def\ed@title{}\let\ed@maketitle\maketitle\renewcommand{\maketitle}{\let\ed@title\@title\ed@maketitle}
\newcommand{\pagestatus}[1]{{\par\centering\small#1\par}}
\newcommand{\zulip}[2][]{{\par\centering\small Discuss on Zulip: \##2, topic
  ``\ifx\relax#1\relax\ed@title\else#1\fi''.\par}\medskip}
% links to other pages of the wiki, anchors, wiki code shown as text
\newcommand{\wikilink}[2]{\ifx\relax#2\relax#1\else#2\fi}
\newcommand{\wikianchor}[1]{\phantomsection\label{#1}}
\newcommand{\wikipage}[1]{}% the first line of a page file of the whole-wiki project (pages/<page>.tex): the page's title
\newcommand{\nowikicode}[1]{\texttt{\frenchspacing#1}}
\newcommand{\nowiki}[1]{#1}
\newcommand{\wikitemplate}[1]{\texttt{\{\{#1\}\}}}
\newcommand{\wikidiv}[2]{\par#2\par}
% \begin{quote}: a sentence set apart, such as the central question of the Overview; centred, in italics
\renewenvironment{quote}{\list{}{\rightmargin\leftmargin}\item\relax\centering\itshape}{\endlist}
% pictures, videos, diagrams: a grey panel with the picture and its caption. The picture (for a video, its poster image)
% is shown when its file, or a .pdf or .png file of the same name, is next to this file, in figures/ or in ../media
\tcbset{ed@media/.style={enhanced,colback=black!4,frame hidden,arc=2pt,left=3mm,right=3mm,top=2.5mm,bottom=2mm,
  before skip=3mm plus 1mm,after skip=3mm plus 1mm}}
\def\ed@poster{}\def\ed@record{}\def\ed@svg{svg}
\define@key{ed@media}{poster}{\def\ed@poster{#1}}\define@key{ed@media}{record}{\def\ed@record{#1}}\define@key{ed@media}{id}{}
\newcommand{\ed@find}[1]{\IfFileExists{#1}{\edef\ed@found{#1}}{\IfFileExists{figures/#1}{\edef\ed@found{figures/#1}}{%
  \IfFileExists{../media/#1}{\edef\ed@found{../media/#1}}{}}}}
\newcommand{\ed@picturefile}[1]{\def\ed@found{}\edef\ed@b{#1}\ifx\ed@b\@empty\else\filename@parse{\ed@b}\edef\ed@base{\filename@base}%
  \edef\ed@ext{\filename@ext}\ifx\ed@ext\ed@svg\else\ed@find{\ed@b}\fi
  \ifx\ed@found\@empty\ed@find{\ed@base.pdf}\fi\ifx\ed@found\@empty\ed@find{\ed@base.png}\fi\fi}
\newcommand{\ed@panel}[2]{\begin{tcolorbox}[ed@media,#1]\ifx\ed@found\@empty\else{\centering\includegraphics[width=\linewidth,
  height=0.35\textheight,keepaspectratio]{\ed@found}\par}\smallskip\fi{\small\textbf{\ifx\ed@a\ed@video Video\else\ifx\ed@a\ed@diagram
  Diagram\else Figure\fi\fi.}\enspace#2\par}{\footnotesize\itshape\ifx\ed@a\ed@video The video plays in the online edition of this page.\fi
  \ifx\ed@record\@empty\else\ifx\ed@a\ed@video\ \fi How it was made: page Numerical methods, record \texttt{\ed@record}.\fi\par}
  \end{tcolorbox}}
% in the text, the panel is a figure (it floats to where it fits); inside a box (a note, a table), it stays in place
\renewcommand{\topfraction}{0.85}\renewcommand{\bottomfraction}{0.6}\renewcommand{\textfraction}{0.1}\renewcommand{\floatpagefraction}{0.8}
\newcommand{\media}[4][]{\par\def\ed@poster{}\def\ed@record{}\setkeys{ed@media}{#1}\def\ed@a{#2}%
  \ifx\ed@a\ed@video\ed@picturefile{\ed@poster}\else\ed@picturefile{#3}\fi
  \ifinner\ed@panel{}{#4}\else\begin{figure}[htbp]\ed@panel{before skip=0pt,after skip=0pt}{#4}\end{figure}\fi}
% tables (\lv{label}{text}; \header{a & b}, \row{a & b}) and code
\newenvironment{levels}[1][]{\par\smallskip\noindent\tabularx{\linewidth}{@{}>{\bfseries\raggedright\arraybackslash}p{0.2\linewidth}X@{}}}{\endtabularx\par\smallskip}
\newcommand{\lv}[2]{#1 & #2\\[3pt]}
\newenvironment{wikitable}[1]{\par\smallskip\noindent\tabularx{\linewidth}{|*{#1}{>{\raggedright\arraybackslash}X|}}\hline}{\endtabularx\par\smallskip}
\newcommand{\header}[1]{#1\\\hline\hline}
\newcommand{\row}[1]{#1\\\hline}
\lstnewenvironment{wikicode}[1][]{\lstset{basicstyle=\ttfamily\footnotesize,breaklines=true,columns=fullflexible,
  keepspaces=true,showstringspaces=false,frame=single}\ifx\relax#1\relax\else\lstset{language=#1}\fi}{}
% the heading of the reference list is the section References of the page
\renewcommand{\bibsection}{}
% characters of the wiki's text that LaTeX does not know
\DeclareUnicodeCharacter{2192}{\ensuremath{\rightarrow}}\DeclareUnicodeCharacter{203A}{\ensuremath{\rangle}}
\DeclareUnicodeCharacter{2713}{\ensuremath{\checkmark}}\DeclareUnicodeCharacter{25CB}{\ensuremath{\circ}}
\DeclareUnicodeCharacter{2248}{\ensuremath{\approx}}\DeclareUnicodeCharacter{2264}{\ensuremath{\leq}}
\DeclareUnicodeCharacter{2265}{\ensuremath{\geq}}\DeclareUnicodeCharacter{2212}{\ensuremath{-}}
\makeatother

\begin{document}
\title{Couette flow}
\author{Wiki \emph{Enhanced dissipation for passive scalars}}
\date{11 November 2026}
\maketitle
\pagestatus{Contributors: Michele Dolce · Status: skeleton}
\zulip{examples}

The Couette flow is the linear shear $u=(y,0)$, considered here in the simplest possible domain for this problem, $(x,y)\in\mathbb{T}\times\mathbb{R}$. The equation under study then reduces to
\[\partial_t f+y\partial_xf=\nu\Delta f, \qquad f|_{t=0}=f^{\rm in}, \qquad (x,y)\in \mathbb{T}\times\mathbb{R}.\]
This is a foundational model in the field, because the same equation governs the vorticity of the two-dimensional Navier–Stokes equations linearized around the Couette flow. In that context Kelvin already solved it explicitly in 1887 (see \wikilink{Literature\#Kelvin1887}{Kelvin (1887, eq. (34)–(35))}). Moreover, the same model, without the diffusion in $x$, was derived by \citet{Kolmogorov1934} when developing his theory of Brownian motion, and the equation is also of Fokker–Planck type. See also \wikilink{History of the problem}{}.

We use this model to illustrate what typical mixing and enhanced dissipation results look like, and how the two are related. The model is simple enough that each proof technique can be carried out in full. They are all presented below on this single example.

\begin{note}[Why this domain?]
The domain $\mathbb{T}\times\mathbb{R}$ is mathematically convenient: there are no boundaries, and one can take the Fourier transform in both variables.

There is also a more geometric motivation. Consider a circular flow in the plane, $u=rV(r)e_\theta$, for which the equation reads $\partial_tf+V(r)\partial_\theta f=\nu\Delta f$. The log-polar coordinates $X=e^s\cos(\theta)$, $Y=e^s\sin(\theta)$ identify $\mathbb{R}^2\setminus\{0\}$ with $(\theta,s)\in\mathbb{T}\times\mathbb{R}$, as shown in the figure below.

\media{image}{log-polar-map.svg}{Left: the plane, with two annuli of the same width and a ray (dashed) carried by a circular flow for some time (solid). Right: the same picture in log-polar coordinates. The annuli become strips, and the circular flow becomes a shear flow.}

Note that this is a conformal change of variables, so the Laplacian only gains a scalar factor, and the equation becomes
\[\partial_t f+V(e^s)\partial_\theta f=\nu e^{-2s}(\partial_\theta^2+\partial_s^2)f,\qquad (\theta,s)\in\mathbb{T}\times\mathbb{R}.\]
If we localize near a radius $r_0=e^{s_0}$ where $V'(r_0)\neq0$, then to leading order, and up to constants, we recover the equation for the Couette flow. In this sense, the Couette flow can be seen as a first local approximation of a radial vortex in $\mathbb{R}^2$, in a region where the angular velocity is strictly monotone. Below is a video showing the effect of the conformal mapping.

\media[poster=log-polar-poster.png,record=num-log-polar]{video}{log-polar.mp4}{Left: a scalar transported and diffused by the circular flow with $V(r)=1/(1+r^2)$. Right: the same solution in log-polar coordinates, where the flow is a shear flow.}
\end{note}

\section{The result}
Before introducing the main theorem, it is useful to observe that, taking the $x$-average, one is left with the standard heat equation on $\mathbb{R}$. Namely,
\[\langle f\rangle_x=\frac{1}{2\pi}\int_{\mathbb{T}}f(x,y)\,dx\quad\Longrightarrow\quad\langle f(t)\rangle_x=e^{\nu t\partial_{yy}}\langle f^{\rm in}\rangle_x.\]
For this part of the solution there is therefore nothing new to say, and it enjoys the standard diffusive estimates. Hence, the quantity to measure is how fast the solution approaches its $x$-average.

\begin{theorem}\label{thm-couette}
Let $\nu\in[0,1)$ and let $f$ be the solution with initial datum $f^{\rm in}\in L^2(\mathbb{T}\times\mathbb{R})$. There exist constants $c,C>0$, independent of $\nu$, such that for all $t\geq0$ the following inequality holds true:
\[\|(f-\langle f\rangle_x)(t)\|_{L^2}\leq Ce^{-c\nu^{1/3}t}\|f^{\rm in}-\langle f^{\rm in}\rangle_x\|_{L^2}.\]
Moreover, if $f^{\rm in}\in H^1(\mathbb{T}\times\mathbb{R})$, the following mixing estimate, uniform in $\nu$, holds true:
\[\|(f-\langle f\rangle_x)(t)\|_{H^{-1}}\leq\frac{C}{1+t}e^{-c\nu^{1/3}t}\|f^{\rm in}-\langle f^{\rm in}\rangle_x\|_{H^1}.\]
\end{theorem}
\begin{note}[The x-average as a projection]
The $x$-average is in fact nothing more than the orthogonal projection in $L^2$ onto the kernel of the transport operator $y\partial_x$: the functions in the kernel are exactly those that do not depend on $x$. When the transport operator has a non-trivial kernel, we cannot expect the elements of the kernel to experience any mixing at all from the transport. The interaction with the dissipation is however much more subtle, and it relies heavily on the commutation properties of this projection with the diffusion operator. In this very simple case everything trivializes: the projection commutes with $\nu\Delta$, so the $x$-average solves the heat equation and $f-\langle f\rangle_x$ solves the same equation as $f$. But the mathematical structure behind the convergence towards the $x$-average lies in the fact that the transport operator has a non-trivial kernel.
\end{note}
\begin{note}[Why negative Sobolev norms measure mixing]
The $L^2$ norm does not see mixing: for $\nu=0$ it is conserved, whatever the flow does. The norm of $H^{-1}$ does, because it measures how much of a function has been sent to high frequencies: in Fourier, it weighs each mode by the inverse of the size of its frequency, so it is small when the function lives at high frequencies, even if its $L^2$ norm is not. On the physical side, negative Sobolev norms scale as powers of a length: the ratio $\|g\|_{H^{-1}}/\|g\|_{L^2}$ can be thought of as the length at which $g$ oscillates. For $\nu=0$, the theorem says that this length decays like $1/t$. See also the note on mixing in the \wikilink{Overview}{}.
\end{note}
\begin{note}[Other function spaces]
The theorem is stated in spaces based on $L^2$. Since the flow is incompressible, the equation does not increase any $L^p$ norm. For $1<p<\infty$, enhanced dissipation in $L^p$ then follows from the $L^2$ estimate by interpolation, with a rate that degenerates as $p\to1$ and as $p\to\infty$. Some techniques give the estimate in $L^p$ directly, without this loss. For the Couette flow, the explicit solution gives the first estimate of the theorem with $L^p$ in place of $L^2$ for every $p\in[1,\infty]$, with constants independent of $p$.
\end{note}
\begin{idea}
The key idea underlying all the methods below is that the transport operator $y\partial_x$ does not commute with the diffusion $\nu\Delta$: indeed, $[\Delta,y\partial_x]=2\partial_x\partial_y$. If the two commuted, the solution would be $e^{-ty\partial_x}e^{\nu t\Delta}f^{\rm in}$ and, since the transport preserves the $L^2$ norm, the flow would have no effect on the decay. We will highlight in each method where this lack of commutation enters.

At first sight, operators that do not commute look like an annoying technical problem. This is a typical instance where a seemingly unfavorable property, looked at in the right way, turns out to be exactly what helps.
\end{idea}
\section{Notation and setting}
Throughout, $\mathbb{T}=\mathbb{R}/2\pi\mathbb{Z}$. For the Fourier transform in $x$ we write
\[f(t,x,y)=\sum_{k\in\mathbb{Z}}f_k(t,y)e^{ikx},\qquad f_k(t,y)=\frac{1}{2\pi}\int_{\mathbb{T}}f(t,x,y)e^{-ikx}\,dx,\]
so that $f_0=\langle f\rangle_x$ and $\|f\|_{L^2}^2=2\pi\sum_{k}\|f_k\|_{L^2(\mathbb{R})}^2$. Note that in this model each $k$-mode evolves independently of the others:
\[\partial_tf_k+ikyf_k=\nu(\partial_{yy}-k^2)f_k,\qquad k\in\mathbb{Z}.\]

The Sobolev norms in the theorem are $\|g\|_{H^1}^2=\|g\|_{L^2}^2+\|\nabla g\|_{L^2}^2$ and its dual,
\[\|g\|_{H^{-1}}=\sup\Big\{\int_{\mathbb{T}\times\mathbb{R}}g\varphi\,dx\,dy\ :\ \|\varphi\|_{H^1}\leq1\Big\}.\]

\section{One example, every technique}
Each proof below is written for this example only, to show the mechanism; the general results, with complete proofs, are in the part \emph{Techniques and general results}.

\begin{note}[The estimates mode by mode]
Since each $k$-mode evolves independently, the estimates of the theorem can be stated for each mode separately, keeping track of the dependence on $k$. There exist constants $c,C>0$, independent of $\nu$ and $k$, such that for all $k\neq0$ and $t\geq0$
\[\|f_k(t)\|_{L^2(\mathbb{R})}\leq Ce^{-c\nu^{1/3}|k|^{2/3}t}\|f_k^{\rm in}\|_{L^2(\mathbb{R})},\]
\[\|f_k(t)\|_{H^{-1}(\mathbb{R})}\leq\frac{C}{1+|k|t}e^{-c\nu^{1/3}|k|^{2/3}t}\|f_k^{\rm in}\|_{H^{1}(\mathbb{R})},\]
where $H^{\pm1}(\mathbb{R})$ are the Sobolev norms in $y$ only. Higher modes are therefore mixed and dissipated faster. The theorem follows by summing over $k\neq0$, since $|k|\geq1$.
\end{note}

\subsection{Explicit solution}
\begin{proof}[Proof by the explicit solution]\label{proof-couette-explicit}
In this proof we take full advantage of the simplicity of the mathematical setting, which is also what allowed Kelvin and Kolmogorov to write down the solution explicitly.

For the equation at mode $k$, we can also take the Fourier transform in $y$, with the notation
\[\hat f_k(t,\eta)=\int_{\mathbb{R}}f_k(t,y)e^{-i\eta y}\,dy,\qquad f_k(t,y)=\frac{1}{2\pi}\int_{\mathbb{R}}\hat f_k(t,\eta)e^{i\eta y}\,d\eta,\]
so that $\|f\|_{L^2}^2=\sum_{k\in\mathbb{Z}}\|\hat f_k\|_{L^2(\mathbb{R})}^2$. By elementary properties of the Fourier transform (multiplication by $y$ becomes $i\partial_\eta$), the equation becomes
\[\partial_t\hat f_k-k\partial_\eta\hat f_k=-\nu(k^2+\eta^2)\hat f_k.\]
This is an advection equation in Fourier, with a damping term. We can then follow the characteristics in $\eta$, which are the lines $\eta+kt=\text{const}$, and deduce that the explicit solution is given by
\[\hat f_k(t,\eta)=\exp\Big(-\nu\int_0^t\big(k^2+(\eta+k\tau)^2\big)\,d\tau\Big)\,\hat f^{\rm in}_k(\eta+kt).\]
\begin{note}[Where the non-commutation enters]
In Fourier, the transport is the translation $\eta\mapsto\eta+kt$ and the diffusion is the multiplication by $-\nu(k^2+\eta^2)$. The lack of commutation is simply reflected in the fact that the damping is not constant along the characteristics:
\[\frac{d}{d\tau}\Big[-\big(k^2+(\eta+k\tau)^2\big)\Big]=-2k(\eta+k\tau),\]
which is exactly the symbol of the commutator $[\Delta,y\partial_x]=2\partial_x\partial_y$ at the frequency $(k,\eta+k\tau)$. If the two operators commuted, the damping would be constant along the characteristics, the integral in the exponent would be trivial, and we would only recover the decay of the heat equation, as in the idea above. Instead, the frequency in $y$ grows linearly in time along the characteristics, so the damping grows like $\nu k^2\tau^2$ and its integral like $\nu k^2t^3$.
\end{note}
The only thing left is therefore to estimate the integral in the exponent, for which we claim that
\[\int_0^t\big(k^2+(\eta+k\tau)^2\big)\,d\tau\geq k^2t+\frac{1}{12}k^2t^3\qquad\text{for all }\eta\in\mathbb{R}.\]
\begin{proof}[Proof of the claim]\label{proof-couette-explicit-claim}
Computing the integral and completing the square in $\eta$,
\[\int_0^t(\eta+k\tau)^2\,d\tau=\eta^2t+\eta kt^2+\frac{1}{3}k^2t^3=t\Big(\eta+\frac{kt}{2}\Big)^2+\frac{1}{12}k^2t^3\geq\frac{1}{12}k^2t^3,\]
with equality only for $\eta=-kt/2$.
\end{proof}
Therefore, using that $f-\langle f\rangle_x=\sum_{k\in\mathbb{Z}\setminus\{0\}}f_ke^{ikx}$, by Plancherel, the fact that $|k|\geq1$ and that $\nu t^3\geq3\nu^{1/3}t-2$, we get
\[\begin{aligned}\|(f-\langle f\rangle_x)(t)\|_{L^2}^2&=\sum_{k\neq0}\int_{\mathbb{R}}e^{-2\nu\int_0^t(k^2+(\eta+k\tau)^2)\,d\tau}\,|\hat f^{\rm in}_k(\eta+kt)|^2\,d\eta\\&\leq C^2e^{-2c\nu^{1/3}t}\,\|f^{\rm in}-\langle f^{\rm in}\rangle_x\|_{L^2}^2.\end{aligned}\]
This proves the desired enhanced dissipation estimate. The constants can be explicitly computed.
\begin{note}[The constants]
Indeed, the claim and $|k|\geq1$ give the more precise bound
\[\|(f-\langle f\rangle_x)(t)\|_{L^2}^2\leq e^{-2\nu t-\frac{1}{6}\nu t^3}\,\|f^{\rm in}-\langle f^{\rm in}\rangle_x\|_{L^2}^2.\]
Moreover, $s^3-3s+2=(s-1)^2(s+2)\geq0$ for $s\geq0$, which for $s=\nu^{1/3}t$ is the inequality $\nu t^3\geq3\nu^{1/3}t-2$. Hence $e^{-\frac{1}{6}\nu t^3}\leq e^{\frac{1}{3}}e^{-\frac{1}{2}\nu^{1/3}t}$, and one can take $c=1/4$ and $C=e^{1/6}$.
\end{note}

For the mixing estimate, one can directly use the characterization of the Sobolev norms in Fourier, namely
\[\|g\|_{H^{s}}^2=\sum_{k\in\mathbb{Z}}\int_{\mathbb{R}}(1+k^2+\eta^2)^{s}\,|\hat g_k(\eta)|^2\,d\eta,\]
which for $s=\pm1$ are the norms of the theorem. Since $|k|\geq1$ and $|k|t\leq|\eta|+|\eta+kt|$, by Cauchy–Schwarz we have
\[\begin{aligned}(1+t)^2\leq(1+|k|t)^2&\leq\big(1+|\eta|+|\eta+kt|\big)^2\leq\big(2+\eta^2\big)\big(2+(\eta+kt)^2\big)\\&\leq\big(1+k^2+\eta^2\big)\big(1+k^2+(\eta+kt)^2\big).\end{aligned}\]
Therefore, bounding the exponential as before,
\[\begin{aligned}\|(f-\langle f\rangle_x)(t)\|_{H^{-1}}^2&=\sum_{k\neq0}\int_{\mathbb{R}}\frac{e^{-2\nu\int_0^t(k^2+(\eta+k\tau)^2)\,d\tau}}{1+k^2+\eta^2}\,|\hat f^{\rm in}_k(\eta+kt)|^2\,d\eta\\&\leq\frac{C^2e^{-2c\nu^{1/3}t}}{(1+t)^2}\sum_{k\neq0}\int_{\mathbb{R}}\big(1+k^2+(\eta+kt)^2\big)|\hat f^{\rm in}_k(\eta+kt)|^2\,d\eta\\&=\frac{C^2e^{-2c\nu^{1/3}t}}{(1+t)^2}\,\|f^{\rm in}-\langle f^{\rm in}\rangle_x\|_{H^1}^2,\end{aligned}\]
which is the mixing estimate.

In fact, the mode-by-mode estimates can be obtained in the same way, by keeping track of $k$.
\begin{note}[The solution in physical space]
For $\nu>0$, inverting the Fourier transform in $y$ in the explicit solution is a Gaussian integral, which gives
\[f_k(t,y)=e^{-\nu k^2t-\frac{1}{12}\nu k^2t^3}\int_{\mathbb{R}}\frac{1}{\sqrt{4\pi\nu t}}\,e^{-\frac{(y-y')^2}{4\nu t}}\,e^{-ikt\frac{y+y'}{2}}f_k^{\rm in}(y')\,dy'.\]
Summing over $k$, by the Poisson summation formula we get
\[f(t,x,y)=\int_{\mathbb{T}\times\mathbb{R}}G_t(x,y,x',y')f^{\rm in}(x',y')\,dx'dy',\]
with
\[\begin{aligned}G_t(x,y,x',y')&=\frac{1}{\sqrt{4\pi\nu t}}\,e^{-\frac{(y-y')^2}{4\nu t}}\sum_{n\in\mathbb{Z}}\frac{1}{\sqrt{2\pi\sigma_t^2}}\exp\Big(-\frac{\big(x-x'-t\frac{y+y'}{2}+2\pi n\big)^2}{2\sigma_t^2}\Big),\\\sigma_t^2&=2\nu t+\frac{1}{6}\nu t^3.\end{aligned}\]
Without the diffusion in $x$, namely with $\sigma_t^2=\frac{1}{6}\nu t^3$, and on the whole plane, namely keeping only the term $n=0$, this is the fundamental solution written by \citet{Kolmogorov1934}. In the variable $x$ the solution spreads over the length $\sqrt{\nu t^3}$, much larger than the diffusive one $\sqrt{\nu t}$. For $t\gg\nu^{-1/3}$ this length exceeds the size of the torus and the kernel is almost independent of $x$: this is the enhanced dissipation seen in physical space.
\end{note}
\end{proof}

\subsection{Spectral methods}
\begin{proof}[Proof by spectral methods]\label{proof-couette-spectral}
\placeholder{The proof.}
\end{proof}
General theory: \wikilink{Spectral methods}{}.

\subsection{Hypocoercivity}
\begin{proof}[Proof by hypocoercivity]\label{proof-couette-hypocoercivity}
\placeholder{The proof.}
\end{proof}
General theory: \wikilink{Hypocoercivity}{}.

\subsection{Stochastic methods}
\begin{proof}[Proof by stochastic methods]\label{proof-couette-stochastic}
\placeholder{The proof.}
\end{proof}
General theory: \wikilink{Stochastic methods}{}.

\subsection{Hörmander theory}
\begin{proof}[Proof by Hörmander theory]\label{proof-couette-hormander}
\placeholder{The proof.}
\end{proof}
General theory: \wikilink{Hörmander theory}{}.

\section{What this example teaches}
\placeholder{What the example shows, and what the general theory keeps of it.}

\section*{References}
%%wiki <!-- refs-begin: Kelvin1887 Kolmogorov1934 -->
\begin{thebibliography}{}
\bibitem[Kolmogorov(1934)]{Kolmogorov1934} A. Kolmogorov, \emph{Zufällige Bewegungen (Zur Theorie der Brownschen Bewegung)}, Annals of Mathematics (2) \textbf{35} (1934), 116–117, \href{https://doi.org/10.2307/1968123}{doi:10.2307/1968123}.
\bibitem[Kelvin(1887)]{Kelvin1887} W. Thomson (Lord Kelvin), \emph{Stability of fluid motion: rectilineal motion of viscous fluid between two parallel planes}, Philosophical Magazine (5) \textbf{24} (1887), 188–196, \href{https://doi.org/10.1080/14786448708628078}{doi:10.1080/14786448708628078}.
\end{thebibliography}
%%wiki <!-- refs-end -->
\end{document}
